\documentclass[10pt,a4paper]{amsart}
\usepackage{amsmath,amssymb,amsthm,graphicx}
% Retained minimal preamble for the Stacks Project chapter topology.tex (Topology).
% The Stacks Project uses one single shared preamble for every chapter; the snapshot stages it
% identically as sources/stacks-topology/preamble.tex, sources/stacks-chow/preamble.tex and
% sources/stacks-etale-cohomology/preamble.tex.  The block below is copied mechanically, line for
% line, from that shared preamble: the Xy-pic setup (lines 14-18) and the native theorem
% declarations together with the mathematical macros (lines 154-217), plus the project's own
% reference/slogan/history display environments (lines 8-12, via the verbatim package).
% Nothing was added, removed or redefined; the wrapper class and the amsmath/amssymb/amsthm
% base packages are supplied by the assembler.
% Chapter: topology.tex -- Topology.
\usepackage[all]{xy}

% We use 2cell for 2-commutative diagrams.
\xyoption{2cell}

% Theorem environments.
%
\theoremstyle{plain}
\newtheorem{theorem}[subsection]{Theorem}
\newtheorem{proposition}[subsection]{Proposition}
\newtheorem{lemma}[subsection]{Lemma}

\theoremstyle{definition}
\newtheorem{definition}[subsection]{Definition}
\newtheorem{example}[subsection]{Example}
\newtheorem{exercise}[subsection]{Exercise}
\newtheorem{situation}[subsection]{Situation}

\theoremstyle{remark}
\newtheorem{remark}[subsection]{Remark}
\newtheorem{remarks}[subsection]{Remarks}

\numberwithin{equation}{subsection}

% Macros
%
\def\lim{\mathop{\mathrm{lim}}\nolimits}
\def\colim{\mathop{\mathrm{colim}}\nolimits}
\def\Spec{\mathop{\mathrm{Spec}}}
\def\Hom{\mathop{\mathrm{Hom}}\nolimits}
\def\Ext{\mathop{\mathrm{Ext}}\nolimits}
\def\SheafHom{\mathop{\mathcal{H}\!\mathit{om}}\nolimits}
\def\SheafExt{\mathop{\mathcal{E}\!\mathit{xt}}\nolimits}
\def\Sch{\mathit{Sch}}
\def\Mor{\mathop{\mathrm{Mor}}\nolimits}
\def\Ob{\mathop{\mathrm{Ob}}\nolimits}
\def\Sh{\mathop{\mathit{Sh}}\nolimits}
\def\NL{\mathop{N\!L}\nolimits}
\def\CH{\mathop{\mathrm{CH}}\nolimits}
\def\proetale{{pro\text{-}\acute{e}tale}}
\def\etale{{\acute{e}tale}}
\def\QCoh{\mathit{QCoh}}
\def\Ker{\mathop{\mathrm{Ker}}}
\def\Im{\mathop{\mathrm{Im}}}
\def\Coker{\mathop{\mathrm{Coker}}}
\def\Coim{\mathop{\mathrm{Coim}}}

% Boxtimes
%
\DeclareMathSymbol{\boxtimes}{\mathbin}{AMSa}{"02}

%
% Macros for moduli stacks/spaces
%
\def\QCohstack{\mathcal{QC}\!\mathit{oh}}
\def\Cohstack{\mathcal{C}\!\mathit{oh}}
\def\Spacesstack{\mathcal{S}\!\mathit{paces}}
\def\Quotfunctor{\mathrm{Quot}}
\def\Hilbfunctor{\mathrm{Hilb}}
\def\Curvesstack{\mathcal{C}\!\mathit{urves}}
\def\Polarizedstack{\mathcal{P}\!\mathit{olarized}}
\def\Complexesstack{\mathcal{C}\!\mathit{omplexes}}
% \Pic is the operator that assigns to X its picard group, usage \Pic(X)
% \Picardstack_{X/B} denotes the Picard stack of X over B
% \Picardfunctor_{X/B} denotes the Picard functor of X over B
\def\Pic{\mathop{\mathrm{Pic}}\nolimits}
\def\Picardstack{\mathcal{P}\!\mathit{ic}}
\def\Picardfunctor{\mathrm{Pic}}
\def\Deformationcategory{\mathcal{D}\!\mathit{ef}}

% --- appended 2026-10-12 (Chief): the project own definitions of the reference/slogan/history
% display environments, copied mechanically from the shared Stacks preamble lines 8-12
% (\usepackage{verbatim} plus the three \newenvironment lines, which the chapter preamble keeps
% verbatim from the project preamble). Nothing added, removed or redefined: these environments
% carry no mathematical content and the project itself typesets them as comments. Without these
% three lines a counted statement whose own text uses them fails with Environment slogan undefined.
\usepackage{verbatim}
\newenvironment{reference}{\comment}{\endcomment}
\newenvironment{slogan}{\comment}{\endcomment}
\newenvironment{history}{\comment}{\endcomment}
\title{Three equivalent descriptions of the topology induced by an injective map}
\author{The Stacks Project Authors}
\date{stacks-project:topology.tex:local-snapshot-20261012; CC BY 4.0}
\begin{document}
\maketitle
\noindent\textit{Editorial note.} All delivered mathematical text is the author's own, copied from the pinned Stacks Project chapter source sources/stacks-topology/pinned-original-source.tex; no statement or proof step was rewritten, reordered or removed.\par\medskip
\noindent\textit{Editorial note.} Preamble provenance: the Stacks Project keeps one single shared preamble for every chapter, staged here as sources/stacks-topology/preamble.tex; this card uses retained-preamble.tex, which copies that shared preamble mechanically (the Xy-pic setup, the native theorem declarations, the mathematical macros and the reference/slogan/history display environments). Nothing was added, removed or redefined; the wrapper class and the base packages are supplied by the assembler.\par\medskip
\noindent\textit{Editorial note.} Bibliography provenance: the chapter carries no chapter-local bibliography (the project keeps one shared bibliography file that is not staged here) and neither the counted statement nor its proof contains any citation, so the card carries no bibliography and no entry was reconstructed.\par\medskip
\noindent\textit{Editorial note.} This card was selected from the staged candidate list for this chapter (candidates-halfA.json / candidates-halfB.json), where it is the candidate labelled lemma-induced; the statement and proof spans are the candidate's own statement\_lines and proof\_lines, checked against the source so that each span opens and closes a real, uncommented environment.\par\medskip
\noindent\textit{Allowed standard prerequisites (editorial).}\begin{itemize}
\item The notion of the topology induced by an injective map of a set into a topological space, and the weakest-topology characterisation. Set up in the author's section on submersive maps, immediately before the counted statement. Fixes the three descriptions whose equivalence the statement asserts.
\end{itemize}\medskip
\begin{lemma}
\label{lemma-induced}
Let $X$ be a topological space. Let $Y$ be a set and let
$f : Y \to X$ be an injective map of sets. The induced
topology on $Y$ is the topology characterized by
each of the following statements:
\begin{enumerate}
\item it is the weakest topology on $Y$ such that $f$ is continuous,
\item the open subsets of $Y$ are $f^{-1}(U)$ for $U \subset X$ open,
\item the closed subsets of $Y$ are the sets $f^{-1}(Z)$ for $Z \subset X$
closed.
\end{enumerate}
\end{lemma}

\begin{proof}
The set $\mathcal{T} = \{f^{-1}(U) | U \subset X \text{ open}\}$ is a
topology on $Y$. Firstly, $\emptyset = f^{-1}(\emptyset)$ and $f^{-1}(X) = Y$.
So $\mathcal{T}$ contains $\emptyset$ and $Y$.

\medskip\noindent
Now let $\{V_i\}_{i \in I}$ be a collection of open subsets where
$V_i \in \mathcal{T}$ and write $V_i = f^{-1}(U_i)$ where $U_i$
is an open subset of $X$, then
$$
\bigcup\nolimits_{i\in I} V_i =
\bigcup\nolimits_{i\in I} f^{-1}(U_i) =
f^{-1}\left(\bigcup\nolimits_{i\in I} U_i\right)
$$
So $\bigcup_{i\in I} V_i \in \mathcal{T}$ as $\bigcup_{i\in I} U_i$
is open in $X$.
Now let $V_1, V_2 \in \mathcal{T}$. We have $U_1, U_2$ open in $X$
such that $V_1 = f^{-1}(U_1)$ and $V_2 = f^{-1}(U_2)$. Then
$$
V_1 \cap V_2 = f^{-1}(U_1) \cap f^{-1}(U_2) = f^{-1}(U_1 \cap U_2)
$$
So $V_1 \cap V_2 \in \mathcal{T}$ because $U_1 \cap U_2$ is open in $X$.

\medskip\noindent
Any topology on $Y$ such that $f$ is continuous contains
$\mathcal{T}$ according to the definition of a continuous map.
Thus $\mathcal{T}$ is indeed
the weakest topology on $Y$ such that $f$ is continuous.
This proves that (1) and (2) are equivalent.

\medskip\noindent
The equivalence of (2) and (3) follows from the equality
$f^{-1}(X \setminus E) = Y \setminus f^{-1}(E)$ for all subsets
$E \subset X$.
\end{proof}

\end{document}
